\section{Concepts in the Generator}\label{sec:concepts}

The probes predict outcomes; sparse autoencoders can tell us what the
model represents about the algorithm it writes. We use public TopK SAEs
for the Qwen3-8B base model (Qwen-Scope, 65,536 features,
$k=100$)~\cite{qwenscope}, which transfer to the post-trained generator
(21--31\% unexplained variance on generated code, over 99\% of the loss
recovered when spliced in). Each feature is summarised per algorithm by
its maximum over the code tokens and matched to the concept labels on
one half of the runs; it is validated on the other half.

Every concept has a feature that detects it (validation AUROC 0.89 to
1.00; step-size adaptation 0.72--0.77), but most of the best features
read keywords (\texttt{mutant}, \texttt{restart}, \texttt{.velocity}). To
separate concept detectors from keyword detectors, we mask every line
that contains a concept keyword (Table~\ref{tab:keyword}). Features for
covariance learning, gradients, particle swarms, surrogates, local
search and adaptive parameters still detect their concept far beyond
what the co-occurring concepts predict (e.g.\ swarm 0.92 against 0.60);
elitism and restart features are keyword detectors, and the crossover
and DE features read the DE family as a whole. At the last prompt token,
the SAE code adds almost nothing to the parent's concepts for predicting
the child's: the concept is not visibly planned before writing.

The concept features say what is implemented, not whether it works:
they add nothing to failure prediction once the code text is known. The
whole SAE code does (+0.074 AUROC), most for \texttt{NameError} (+0.19)
and \texttt{AttributeError} (+0.18). One layer-21 feature fires on about
half of all undefined names and on 0.3\% of valid ones, for instance the
objective function used outside its scope. Steering on it raises the
rate of undefined names from 5\% to 30\%, so it is causal for this
error class.

Features that \emph{read} a concept are candidate directions to
\emph{write} it, which we test next.
